\clearpage
\appendix
\section{问题一及共用核心程序}

第一至第四问共用径向有限体积空间离散。文件\texttt{1.txt}首先给出物性函数、表层加密网格、守恒通量、隐式时间推进、临界时刻定位和Richardson误差估计，随后给出问题一的常物性温度场与非线性水分扩散实现。附件读取、结果写入、日志和绘图程序从略，完整可执行程序见支撑材料。

\VerbatimInput[
  fontsize=\scriptsize,
  numbers=left,
  numbersep=8pt,
  xleftmargin=1.8em,
  breaklines=true,
  breakanywhere=false,
  breaksymbolleft={},
  breaksymbolright={},
  frame=single,
  framesep=2mm
]{code/1.txt}

\section{问题二完整程序的关键实现}

以下摘录给出表面加密网格、控制体权重、Jacobian稀疏结构、变物性耦合方程、自适应BDF求解及空间收敛计算。

\VerbatimInput[
  fontsize=\scriptsize,numbers=left,numbersep=8pt,xleftmargin=1.8em,
  breaklines=true,breakanywhere=false,breaksymbolleft={},breaksymbolright={},
  frame=single,framesep=2mm
]{code/2.txt}

\section{问题三完整程序的关键实现}

以下摘录给出后向Euler启动、BDF2时间推进、Picard耦合迭代、全域阈值事件判断及临界时刻定位。

\VerbatimInput[
  fontsize=\scriptsize,numbers=left,numbersep=8pt,xleftmargin=1.8em,
  breaklines=true,breakanywhere=false,breaksymbolleft={},breaksymbolright={},
  frame=single,framesep=2mm
]{code/3.txt}

\section{问题四完整程序的关键实现}

以下摘录给出时变半径下的固定域变换、移动边界热湿方程、连续临界时刻求解、干质量后验检验及三级网格收敛分析。

\VerbatimInput[
  fontsize=\scriptsize,numbers=left,numbersep=8pt,xleftmargin=1.8em,
  breaklines=true,breakanywhere=false,breaksymbolleft={},breaksymbolright={},
  frame=single,framesep=2mm
]{code/4.txt}
